% !PRINTED 137
\weberpage{137}

Ist \(\omega\) der Modul der \(\vartheta\)-Funktionen, so gehören, wie sich nach (1) aus
dem Verschwinden der vier \(\vartheta\)-Funktionen ergibt, die folgenden Werte
zusammen:
\weq{11}{\begin{aligned}
u &= \tfrac{1}{2}\omega, & \xi : \eta &= \xi_1 : \eta_1,\\
u &= 0, & \xi : \eta &= \xi_2 : \eta_2,\\
u &= \tfrac{1}{2}, & \xi : \eta &= \xi_3 : \eta_3,\\
u &= \frac{1+\omega}{2}, & \xi : \eta &= \xi_4 : \eta_4.
\end{aligned}}

\begin{center}
\textbf{\S\,42. \gesperrt{Jacobis elliptische Funktionen.}}
\end{center}

Da man die Variablen \(\xi\), \(\eta\) mittels einer linearen Substitution, in der
vier Koeffizienten disponibel sind, durch zwei neue Variable ersetzen
kann, so kann man vier von der Größe \(\xi_i\), \(\eta_i\) oder drei von ihren
Verhältnissen beliebige Werte erteilen, ohne die Allgemeinheit zu
beeinträchtigen.

Wir wollen setzen
\weq{1}{\xi_1 = 0,\qquad \eta_2 = 0,\qquad \xi_3 = \eta_3,}
und führen noch \(\varkappa^2\), \(\varkappa'^2\), \(\zeta\) durch die Gleichungen an:
\weq{2}{\xi_4 = \varkappa^2\eta_4,\qquad \varkappa'^2 = 1 - \varkappa^2,\qquad \frac{\eta}{\xi} = \zeta.}

Aus \S\,41, (3), (5), (6) ergibt sich dann
\weq{3}{\begin{aligned}
\frac{\xi_2}{\eta_1} &= -\frac{\vartheta^2_{10}}{\vartheta^2_{00}}, &
\frac{\eta_3}{\eta_1} &= \frac{\vartheta^2_{10}}{\vartheta^2_{01}}, &
\frac{\eta_4}{\eta_1} &= \frac{\vartheta^2_{00}}{\vartheta^2_{01}},\\
\sqrt{\varkappa} &= \frac{\vartheta_{10}}{\vartheta_{00}}, &
\sqrt{\varkappa'} &= \frac{\vartheta_{01}}{\vartheta_{00}}, & &
\end{aligned}}
und aus (1), (9), \S\,41 findet sich:
\weq{4}{\begin{aligned}
\frac{\vartheta_{00}}{\vartheta_{10}}\,\frac{\vartheta_{11}(u)}{\vartheta_{01}(u)} &= \sqrt{\zeta},\\
\frac{\vartheta_{01}}{\vartheta_{10}}\,\frac{\vartheta_{10}(u)}{\vartheta_{01}(u)} &= \sqrt{1-\zeta},\\
\frac{\vartheta_{01}}{\vartheta_{00}}\,\frac{\vartheta_{00}(u)}{\vartheta_{01}(u)} &= \sqrt{1-\varkappa^2\zeta}.
\end{aligned}}
\weq{5}{2\pi\vartheta^2_{00}u = \int_0^{\zeta} \frac{d\zeta}{\sqrt{\zeta(1-\zeta)(1-\varkappa^2\zeta)}},}
worin der Integrationsweg und die Bedeutung der Wurzelzeichen durch
die Formeln (4) selbst bestimmt ist, wenn der Übergang
%CONT>>
